\documentclass[10pt,a4paper]{amsart}
\usepackage[margin=19mm]{geometry}
\usepackage{amsmath,amssymb,mathtools}
\usepackage[scr=rsfs,cal=euler]{mathalfa}
\usepackage{xfrac}
\usepackage[unicode,hidelinks,bookmarks=false]{hyperref}
\newcommand{\ie}{i.e.\ }
\newcommand{\cf}{cf.\ }
\newcommand{\resp}{respectively\ }
\newcommand{\Subsection}{\subsection}
\providecommand{\nobreakdash}{\nobreak\hbox{-}}
\setlength{\emergencystretch}{2em}
\multlinegap0pt
\usepackage{mathtools}
\usepackage{amssymb}
\usepackage{bm}
\usepackage{xcolor}
\usepackage{float}
\newtheorem{lemma}{Lemma}
\newcommand{\E}{\mathbb{E}}
\newcommand{\FF}{\mathcal{F}}
\newcommand{\PP}{\mathcal{P}}
\newcommand{\al}{\alpha}
\newcommand{\kk}{\kappa}
\newcommand{\ol}{\overline}
\newcommand{\rr}{\rho}
\newcommand{\wt}{\widetilde}
\title{Computation of a separatrix map and a normally hyperbolic invariant lamination for the RP3BP}
\author{Marcel Guardia, Vadim Kaloshin, Pau Mart\'in, Pablo Roldan}
\date{Source: arXiv:2603.19893v2 (CC BY 4.0)}
\begin{document}
\maketitle

\noindent\emph{Source and licence.} Computation of a separatrix map and a normally hyperbolic invariant lamination for the RP3BP. Version arXiv:2603.19893v2 (CC BY 4.0), distributed by arXiv under the Creative Commons Attribution 4.0 International licence, http://creativecommons.org/licenses/by/4.0/, as recorded in the arXiv licence metadata for this exact version and shown on the abstract page https://arxiv.org/abs/2603.19893v2. Full licence notice: \texttt{licenses/arxiv-2603.19893-CC-BY-4.0.txt}.\par\smallskip

\noindent\textit{Editorial scope.} Counted unit: the article's lemma lemma:DecayConditionned (source0.tex lines 4541--4553). The article's complete original proof of it is delivered with it (source0.tex lines 4651--4759, one \texttt{proof} environment of 109 lines, which the article places away from the statement). No mathematics is written, repaired or re-derived: the statement and the proof are the article's own lines, in the article's own notation, and no retained line is substituted or re-flowed.  Retained uncounted as the source-local context the proof needs: lines 4356--4359; lines 4407--4411.  Nothing was omitted from any retained span, so the package carries no omission record.  No retained line contains a citation, so the wrapper carries no bibliography and no bibliography entry.  No retained line contains a figure environment, an image-inclusion command or drawing code, so no image is delivered and the package declares no asset dependency.  Editorial formatting only, in addition: an AMS \texttt{amsart} preamble that restates the article's own declarations and macros used by the retained lines, namely the environments \texttt{lemma} on separate counters, together with the standard packages those lines need.  The retained environments print in this document as Lemma 1 in order of appearance, whereas the article prints the same items with the numbering of its own full document.  The complete native source of the article as distributed in the arXiv e-print for this exact version is bundled unchanged as \texttt{sources/196808/source0.tex}.
\begin{equation}\label{def:DecayCorrExpectheta}
 \left|\E_\omega \phi\circ\sigma_h^i\psi^T\sigma_h^j\right|\leq 
C\|\phi\|_\theta \|\psi\|_\theta r^{i-j}.
\end{equation}

\begin{align}\label{def:xieta}
 \xi_j&=\sum_{m=jp+jq+1}^{(j+1)p+jq}X_m,\qquad \text{for}\qquad 0\leq j\leq k-1\\
\eta_j&=\sum_{m=(j+1)p+jq+1}^{(j+1)p+(j+1)q}X_m,\qquad \text{for}\qquad 0\leq j\leq k-1\\
\eta_k&=\sum_{m=kp+kq+1}^{n}X_m.
\end{align}

\begin{lemma}\label{lemma:DecayConditionned}
% By \eqref{def:DecayCorrExpectheta}, it is easy to check that 
\[
 \begin{split}
 \left|\E_\omega\left(\left.
 \xi_j\right|\FF_{j-1}\right)-\E_\omega\left(
 \xi_j\right)\right|\lesssim r^{q}\\
\left| \E_\omega\left(\left.
\xi_j^2\right|\FF_{j-1}\right)-\E_\omega\left(
 \xi_j^2\right)\right|\lesssim r^{q}
  \end{split}
\]
\end{lemma}

\begin{proof}[Proof of Lemma \ref{lemma:DecayConditionned}]
We first obtain the first estimate. Note that by hypothesis, $\E_\omega(\xi_j)=0$.
Then, note that, by \eqref{def:xieta},
\[
\E_\omega(\xi_j|\FF_{j-1})=\sum_{m=jp+jq+1}^{(j+1)p+jq}\E_\omega(X_m|\FF_{j-1}).
 \]
The random variables $X_m$ only have harmonics $\pm 1$. Conditioning by $\FF_{j-1}$ corresponds to fixing the values of the symbols $\{\omega_0,\ldots, \omega_{jp+(j-1)q}\}$. Therefore, we are in the same setting as in \eqref{def:DecayCorrExpectheta}.

Now we prove the second estimate in Lemma \ref{lemma:DecayConditionned}. Note that 
\begin{equation}\label{def:xideg2}
\xi_j^2=\sum_{m=jp+jq+1}^{(j+1)p+jq}X_m^2+2\sum_{m,k=jp+jq+1, m<k}^{(j+1)p+jq}X_mX_k
\end{equation}
Each of these terms have harmonics $0$ and $\pm 2$. Let us start by the harmonic $0$. It contains terms of the form 
\[
\PP= P(\sigma^k\omega)\ol P(\sigma^m\omega)e^{i\sum_{i=m}^{k-1}\beta(\sigma^i\omega)}.
\]
% When $n=m$, we have the term
% \[
% \E_\omega(| P(\sigma^n\omega)|^2|\FF_{j-1})-\E_\omega(| P(\sigma^n\omega)|^2)
% \]
% Now, we write $\E_\omega=\E_{\omega_\leq}\E_{\omega_>}$, where $\E_{\omega_>}$ refers to the expectation with respect to the symbols $\{\omega_k\}_{k\geq jp+(j-1)q+1}$ and $\E_{\omega_\leq}\E_{\omega_>}$ refers to the expectation with respect to $(\omega_0,\ldots, \omega_{jp+(j-1)q})$. Then, we can write 
% \[
%  \begin{split}
%   \E_\omega(| P(\sigma^n\omega)|^2|\FF_{j-1})-\E_\omega(| P(\sigma^n\omega)|^2)&=  \E_{\omega,>}(| P(\sigma^n\omega)|^2)-\E_{\omega,>}\E_{\omega,\leq}(| P(\sigma^n\omega)|^2)\\
%   &=\E_{\omega,>}\left[| P(\sigma^n\omega)|^2)-\E_{\omega,\leq}(| P(\sigma^n\omega)|^2)\right]
%  \end{split}
% \]
% Then, since $P$ is H\"older with respect to $\omega$, one has that there exists $r>0$ such that 
% \[
%  \left|| P(\sigma^n\omega)|^2)-\E_{\omega,\leq}(| P(\sigma^n\omega)|^2)\right|\lesssim r^{\log^2n},
% \]
% which implies 
% \[
% \left|\E_\omega(| P(\sigma^n\omega)|^2|\FF_{j-1})-\E_\omega(| P(\sigma^n\omega)|^2)\right|\lesssim r^{\log^2n}.
% \]
% For $m\neq n$ we proceed as 
We need to estimate the term
\[
\E_\omega\left(\left. \PP\right|\FF_{j-1}\right)-\E_\omega\left( \PP\right)
\]
Now, we write $\E_\omega=\E_{\omega_\leq}\E_{\omega_>}$, where $\E_{\omega_>}$ refers to the expectation with respect to the symbols $\{\omega_k\}_{k\geq jp+(j-1)q+1}$ and $\E_{\omega_\leq}$ refers to the expectation with respect to $\{\omega_0,\ldots, \omega_{jp+(j-1)q}\}$. Then, we can write 
\[
 \begin{split}
\E_\omega\left(\left. \PP\right|\FF_{j-1}\right)-\E_\omega\left( \PP\right)
&=  \E_{\omega,>}(\PP)-\E_{\omega,>}\E_{\omega,\leq}(\PP)\\
  &=\E_{\omega,>}\left[\PP-\E_{\omega,\leq}(\PP)\right]
 \end{split}
\]
Then, since $P$ is H\"older with respect to $\omega$, one has that there exists $r>0$ such that 
\[
 \left|\PP-\E_{\omega,\leq}(\PP)\right|\lesssim r^{\log^2n},
\]
which implies 
\[
\left|\E_\omega(\PP|\FF_{j-1})-\E_\omega(\PP)\right|\lesssim r^{\log^2n}.
\]
Now we consider the harmonic 2 in \eqref{def:xideg2} (the harmonic -2 can be estimated analogously). It contains terms of the form 
\[
\PP= P(\sigma^k\omega)\ol P(\sigma^m\omega)e^{i\sum_{i=m}^{k-1}\beta(\sigma^i\omega)}e^{2i\sum_{i=0}^{m-1}\beta(\sigma^i\omega)}.
\]
We use again the decay of correlations in \eqref{def:DecayCorrExpectheta}. We consider two cases depending $|k-m|$. 

Note that we have two parameters: $\rr$ given by the decay of correlations in \eqref{def:DecayCorrExpectheta} and $\theta$ coming from the $\theta$-H\"older dependence of $P$ and $\beta$ on $\omega$. We choose $\kk$ such that 
\[
 \theta^{-\kk}\rr\leq \wt \rr<1.
\]
Then, we consider first $k-m\leq \kk q$ and then $k-m\geq \kk q$. 

For the first case, we apply the decay of correlations  in \eqref{def:DecayCorrExpectheta} to the functions 
\[
\phi(\al,\omega)=Q(\omega)e^{i2\al}+ Q(\omega)e^{-i2\al},\quad \psi(\al,\omega)=e^{i2\al}+e^{-i2\al}
\]
where
\[
 Q(\omega)=P(\sigma^{k-m}\omega)P(\omega)e^{2i\sum_{i=0}^{k-m-1}\beta(\sigma^i\omega)}
\]
% where
% \[
%  Q(\omega)=P(\omega)e^{2i\sum_{i=0}^{n-m-1}\beta(\sigma^i\omega)}
% \]
Then,
\[
 |\E_\omega(\PP|\FF_{j-1})|\leq \|Q\|_\theta \rr^q
\]
Since $\|Q\|_\theta\leq \theta^{-|k-m|}$ and $k-m\leq \kk q$, we have
\[
 |\E_\omega(\PP|\FF_{j-1})|\leq \theta^{-\kk q}\rr^q \leq \wt\rr^{\log^2n}.
\]
For the second case, $k-m\geq \kk q$, we further condition the expectation to 
\[
\E_\omega(\PP|\FF_{j-1})=\E_\omega\left(\E_\omega(\PP|\{\omega_k\}_{k\leq m-1})|\FF_{j-1}\right).
\]
Note that,
\[
 \E_\omega(\PP|\{\omega_k\}_{k\leq m-1})= e^{i\sum_{i=0}^{m-1}\beta(\sigma^i\omega)}\E_\omega\left(\left. P(\sigma^k\omega)\ol P(\sigma^m\omega)e^{2i\sum_{i=m}^{k-1}\beta(\sigma^i\omega)}\right|\{\omega_k\}_{k\leq m-1}\right)
\]
Now, for any fixed values for $\{\omega_k\}_{k\leq m-1}$, one can apply \eqref{def:DecayCorrExpectheta} to obtain 
\[
 \left|\E_\omega\left(\left. P(\sigma^k\omega)\ol P(\sigma^m\omega)e^{2i\sum_{i=m}^{k-1}\beta(\sigma^i\omega)}\right|\{\omega_k\}_{k\leq m-1}\right)\right|\leq r^{k-m}.
\]
Since by hypothesis $k-m\geq \kk q$, we obtain 
\[
|\E_\omega(\PP|\FF_{j-1})|\leq \rr^{\log^2 n}.
\]
Note that one can obtain analogos estiamtes for $\E_\omega(\PP)$ just writing
\[
\E_\omega\left( \E_\omega(\PP|\FF_{j-1})\right).
\]
\end{proof}

\end{document}
